%% kapa: exam=1 round=35 year=2024 period=1 form=A subject=civil_law subject_ko=민법 seq=8 number=8
%% source: Q-Net 감정평가사 기출문제 (cst00309) artlSeq=5213229 file=raw/1008/5213229/1교시 문제지 원본.pdf
%% conversion: kordoc 4.24.51 -> tools/gen1.py (auto); review=proofread-sonnet; flags=none
\begin{question}{8}
권리의 객체에 관한 설명으로 옳지 않은 것은? (다툼이 있으면 판례에 따름)
\begin{choices}
  \item 토지의 개수는 「공간정보의 구축 및 관리 등에 관한 법률」에 의한 지적공부상 토지의 필수(筆數)를 표준으로 결정된다.
  \item 1필의 토지의 일부가 「공간정보의 구축 및 관리 등에 관한 법률」상 분할절차 없이 분필등기가 된 경우, 그 분필등기가 표상하는 부분에 대한 등기부취득시효가 인정될 수 있다.
  \item 주물에 대한 점유취득시효의 효력은 점유하지 않은 종물에 미치지 않는다.
  \item 주물의 상용에 제공된 X동산이 타인 소유이더라도 주물에 대한 경매의 매수인이 선의취득 요건을 구비하는 경우, 그 매수인은 X의 소유권을 취득할 수 있다.
  \item 명인방법을 갖춘 미분리과실은 독립한 물건으로서 거래의 객체가 될 수 있다.
\end{choices}
\end{question}
\answer{2}
